\section{Diseño e Implementación}
\subsection{Hardware del sistema}
A continuación se muestra un diagrama de bloques del Hardware del equipo y se explicarán las conexiones para cada uno de los componentes. 
\subsubsection{Display LCD 16x2}
\begin{table}[H]
\centering
\begin{tabular}{|c|c|}
\hline
Nucleo & Display\\ \hline
5V& VSS \\ \hline
5V & VDD \\ \hline
B12 & RS \\ \hline
GND & RW \\ \hline
A11 & E \\ \hline
- & D0 \\ \hline
- & D1 \\ \hline
- & D2 \\ \hline
- & D3 \\ \hline
PB5 & D4 \\ \hline
PB4 & D5 \\ \hline
PB13 & D6 \\ \hline
PB14 & D7 \\ \hline
GND & K \\ \hline
\end{tabular}
\caption{Conexiones del Display}
\end{table}
El pin A se conecta a 5V a través de una resistencia de $1k\ohm$ mientras que el pin Vo se conecta a 5V y GND a través de un potenciometro de $100k\ohm$.
\subsubsection{Sensor MAX30102}
\begin{table}[H]
\centering
\begin{tabular}{|c|c|}
\hline
Nucleo & Sensor MAX30102 \\ \hline
GND & GND \\ \hline
3.3V & VCC \\ \hline
PB6 & SCL \\ \hline
PB7 & SDA \\ \hline
\end{tabular}
\caption{Conexiones del sensor MAX30102}
\end{table}

\subsubsection{Módulo bluetooth HM-10}
\begin{table}[H]
\centering
\begin{tabular}{|c|c|}
\hline
Nucleo & HM-10 \\ \hline
GND & GND \\ \hline
5V & VCC \\ \hline
PA3 & TXD \\ \hline
PA2 & RXD \\ \hline
\end{tabular}
\caption{Conexiones del módulo bluetooth HM-10}
\end{table}

\subsubsection{Buzzer}
\begin{table}[H]
\centering
\begin{tabular}{|c|c|}
\hline
Nucleo & Buzzer \\ \hline
GND & - \\ \hline
PC7 & + \\ \hline

\end{tabular}
\caption{Conexiones del sensor MAX30102}
\end{table}

\subsubsection{Botones}
\begin{table}[H]
\centering
\begin{tabular}{|c|c|}
\hline
Nucleo & BOTON \\ \hline
PC10 & SWITCH-MODE \\ \hline
PC11 & ALARM-OFF \\ \hline
\end{tabular}
\caption{Conexiones de los botones}
\end{table}
Siendo SWITCH-MODE y ALARM-OFF los botones encargados de cambiar el modo de uso y apagar la alarma respectivamente, ambos configurados en modo pull-up y conectados a GND del otro extremo.

\subsubsection{LEDs}
Para polarizar los LEDs con una fuente de 3.3V, las resistencias limitadoras se calculan mediante la siguiente fórmula: 
\begin{equation}
    R = \frac{V_{CC}- V_f}{I_f}
\end{equation}

\begin{itemize}
    \item $V_{CC}$ la tensión de alimentación (en este caso 3.3V).
    \item $I_f$ la corriente deseada por el LED(normalmente se utiliza entre 5ma y 10ma).
    \item $V_f$ la tensión directa del LED.
\end{itemize}
Dado que se utilizan LEDs de diferentes colores el valor de $V_f$ dependerá de la caída de tensión típica según el color: 
\begin{itemize}
    \item LED rojo $\longrightarrow$ indica estado de peligro junto con la alarma. 
    \item LED azul $\longrightarrow$ indica que la conexión bluetooth a sido exitosa.
    \item LED verde $\longrightarrow$ modo adulto. 
    \item LED amarillo $\longrightarrow$ modo niño.
\end{itemize}

\begin{table}[H]
\centering
\begin{tabular}{|c|c|}
\hline
Color LED & $V_f$ típico \\ \hline
Rojo & 2V \\ \hline
Amarillo & 2.1V \\ \hline
Verde & 2.2V \\ \hline
Azul & 3.1V \\ \hline
\end{tabular}
\caption{Caída de tensión típica según el color}
\end{table}
Se tomó como valor deseado de corriente $I_f = 10mA$.
Los valores de las resistencias utilizadas para los LEDs amarillo, verde y rojo son de $R = 220 \ohm$ y $R= 30\ohm$ para el LED azul, conectados como se observa en la figura \ref{fig:LEDs}. 

\begin{figure}[H]
    \centering
    \begin{circuitikz}
        \draw (-1,1) node{3,3V} (-.5,1) node[circ]{} -- (0,1)
        to [led] (2,1)
        to [R] (4,1) -- (4.5,1)
        -- ++(0,-1) node[tlground]{};
    \end{circuitikz}
    \caption{Conexión de los LEDs.}
    \label{fig:LEDs}
\end{figure}


\begin{table}[H]
    \centering
    \begin{tabular}{|c|c|} \hline
        Nucleo & LED \\ \hline
        PB10 & LED-Alarm \\ \hline
        PB11 & LED-Bluetooth \\ \hline
        PB2 & LED-ADULT \\ \hline
        PB1 & LED-KID \\ \hline
    \end{tabular}
    \caption{Conexión de los LEDs}
\end{table}

\subsection{Costo de componentes}

\begin{table}[H]
    \centering
    \begin{tabular}{|c|c|} \hline
        Componente & Costo\\ \hline
        MAX30102& \\ \hline
        HM10&  \\ \hline
        Display&  \\ \hline
        Resistencias& \\ \hline
        LEDs& \\ \hline
        Buzzer&  \\ \hline
        Botones&  \\ \hline
    \end{tabular}
    \caption{Costos de componentes}
\end{table}

\subsection{Firmware}
%
% El firmware del dispositivo se programó en lenguaje C.
% La división de responsabilidades se realizó en los siguientes módulos.

\subsubsection{Sensor}\label{sec:sensor} % (fold)
Este módulo representa un sensor que se comunica con un oxímetro de pulso (\textit{MAX30102}) y procesa la señal para obtener los datos en la siguiente estructura de datos.%~\ref{data_struct}.

%formatear mas lindo con colorcitos y eso
\begin{verbatim}
    typedef struct {
        uint8_t heart_rate;
        uint8_t respiratory_rate;
        uint8_t spo2;
        uint8_t apnea;
        uint32_t timestamp;
    } task_sensor_results_dta_t;
\end{verbatim}

La comunicación se realiza via $I^2C$ y en cada iteración se acceden los datos al buffer FIFO del MAX30102 para reconstruir la señal.
Una vez procesada, cada \texttt{DEL\_SEN\_MAX\_TICKS} se accede a la interfaz del sistema y se actualiza la estructura de datos, además de activar una variable booleana indicando que el valor no ha sido leído.

% subsubsection Sensor (end)

\subsubsection{Actuador}\label{sec:actuador} % (fold)
Consiste en los indicadores al usuario.
Consta de 4 estados.
Pueden estar prendidos, apagados, parpadeando o emtiendo un pulso.
Éstos últimos permiten ajustar la frecuencia con la que un led parpadea, o la frecuencia de un buzzer.
% subsubsection Actuador (end)

\subsubsection{Botones}\label{sec:botones} % (fold)
Modela el comportamiento de un pulsador mecanico.
Permite filrar por software el ruido durante los cambios de estado del botón (\textit{debouncing}).
Aunque el botón es mecánico, se puede modelar otras entradas digitales con ruido de la misma manera.
% subsubsection Botones (end)


\subsubsection{Sistema}\label{sec:sistema} % (fold)
Es la tarea principal encargada de recibir las entradas desde los botones y sensores y producir los efectos correspondientes en los otros módulos.
Como es relevante no perder ninguna entrada, cuenta con una interfaz que gestiona una cola de eventos.
Ésta, además, es donde el sensor guarda los datos que toma cada cierta cantidad de ticks.
En cada iteración, el sistema pregunta a la cola si hay algun evento disponible y si hay algun dato para enviar.
En caso de haber un dato disponible del sensor, lo envía por UART al HM10, que se encarga de transmitirlo por bluetooth, y se muestra por el display

% subsubsection Sistema (end)

\subsection{Diseño}\label{sub:diseño} % (fold)

Se utilizó \textit{KiCAD} durante la etapa de diseño, para preparar la pcb del dispositivo.

\begin{figure}[H]
    \centering
    \colorbox{kicadpcbbg}{\includegraphics[width=.55 \textwidth]{tp3.png}}
    \caption{Diseño de la PCB en KiCAD.}
\end{figure}


La disposición de la placa fue pensada para soldarse en una placa experimental.
Se representaron los cables como pistas y se tuvo en cuenta el espaciado entre orificios.

\begin{figure}[H]
    \centering
    \includegraphics[height=.8 \textwidth]{render_side.png}
    \hspace{4pt}
    \includegraphics[height=.8 \textwidth]{render.png}
    \caption{Previsualización 3D de la placa.}
    \label{fig:}
\end{figure}

Se cuenta con espacio suficiente para agregar una batería para hacer al dispositivo portable.
Considerando que solo se precisa el integrado y el uso de una pantalla con menor tamaño y similar o mejor resolución queda margen para un modelo más compacto. 

{\color{red}en conclusiones tal vez?}

% subsection Diseño (end)

\subsection{Ensamblaje}\label{sub:ensamblaje} % (fold)
% subsection Ensamblaje (end)
